\chapter{IA, Aprendizagem e Modelos: Começar do Zero}
\label{ch:ia}

\ficcao A Ana recebe um email da biblioteca do bairro. A reserva do livro ficou registada. Antes de sair de casa, pergunta se já o pode levantar. No balcão, o registo mostra «pedido recebido». A colega que organiza as reservas ainda está a conferir os exemplares. Quem redige a resposta tem uma frase muito fácil de escrever, «O livro está à sua espera», e uma informação que lhe falta para a poder usar.

É por aí que vamos começar. O que queremos saber? Que informação temos? Como confirmamos a resposta? Podes seguir o caso com papel e lápis.

\section{Uma pergunta pode esconder três trabalhos}

Se quisermos contar quantas reservas entraram hoje, basta contar os registos desse dia. Se quisermos separar centenas de mensagens em «reservas», «devoluções» e «horários», podemos usar um sistema que reconheça formas diferentes de fazer a mesma pergunta. Se quisermos confirmar o estado da reserva da Ana, precisamos do registo correspondente e de saber o que significa o estado nele escrito.

No primeiro caso, a regra é uma contagem. No segundo, a tarefa é de \emph{classificação}, isto é, escolher uma categoria para cada mensagem. No terceiro, é uma consulta. Redigir a resposta é ainda outra operação. Separar estes trabalhos ajuda a escolher a ferramenta e a detetar onde falta informação.

\section{O que a máquina aprende}

A inteligência artificial (IA) reúne abordagens para tarefas como reconhecer padrões, procurar uma solução ou escolher uma ação. Num programa de regras, alguém escreve as condições. Por exemplo, quando a reserva tem o estado «pronto para levantamento», a aplicação pode preparar um aviso. Na \emph{aprendizagem automática}, ou \emph{machine learning}, usamos dados para ajustar um modelo. A partir de mensagens já classificadas, o modelo pode aprender padrões que o ajudam a tratar mensagens novas.

Chama-se \emph{generalização} a capacidade de funcionar em casos novos. Acertar nas mensagens usadas para aprender é uma parte do trabalho. Para saber se o modelo também reconhece pedidos escritos de outras formas, testa-o com mensagens que ficaram de fora. O capítulo~\ref{ch:ml} abre essa experiência com números pequenos.

Uma rede neuronal organiza operações em etapas chamadas \emph{camadas}. Os \emph{parâmetros} são os números que controlam essas operações. Pensa na regra de uma conta: se cada quilo acrescenta dois minutos, o dois controla o resultado. Numa rede, o treino ajusta muitos números deste género. Na \emph{aprendizagem profunda}, ou \emph{deep learning}, várias camadas transformam a entrada em representações que ajudam na tarefa~\citep{goodfellow2016}.

\begin{keyideabox}{Treino, modelo e aplicação}
O \textbf{algoritmo de treino} é o procedimento que ajusta os parâmetros. O \textbf{modelo treinado} usa os valores obtidos. A \textbf{aplicação} organiza o pedido, os dados, o modelo e as ferramentas. A mesma aplicação pode procurar um registo, calcular uma data por regras e pedir ao modelo que explique o resultado.
\end{keyideabox}

\section{Gerar uma resposta e consultar um registo}

Um modelo \emph{generativo} pode produzir texto, imagem, som ou outros tipos de saída, conforme o seu desenho. Um modelo de linguagem aprende regularidades de texto. Nos modelos que geram texto por continuação, a entrada é dividida em \emph{tokens}, unidades que podem ser palavras, fragmentos ou sinais. O modelo atribui probabilidades a unidades que podem aparecer a seguir. Um LLM, do inglês \emph{large language model}, é um modelo de linguagem de grande escala.

A frase «Pode levantar o livro hoje» é uma continuação plausível de um pedido sobre uma reserva. Para a usar, o balcão precisa de confirmar que aquele exemplar foi separado para aquela pessoa. Se a aplicação tiver acesso ao registo, pode consultá-lo antes de gerar o texto. Se a informação ficar por confirmar, a resposta pode indicar o que já se sabe e o que falta.

\section{O que muda quando escrevemos um pedido}

No \emph{treino}, ajustamos os parâmetros do modelo. Na \emph{inferência}, usamos os parâmetros para produzir uma saída. Quando colocas a regra da biblioteca na conversa, mudas a informação disponível nessa utilização. O capítulo~\ref{ch:contexto} mostra como o resultado pode mudar com exemplos e informação adicional. O ajuste posterior dos parâmetros, ou \emph{fine-tuning}, é tratado no capítulo~\ref{ch:adaptacao}.

O \emph{prompt} é o pedido e as instruções que fornecemos ao modelo. O \emph{contexto} é o conjunto de informação que o sistema lhe disponibiliza naquele momento. Pode incluir o pedido, mensagens anteriores, um documento ou o resultado de uma consulta. Pensa numa reunião. Uma coisa é responder com a agenda à frente. Outra é ter também o relatório atualizado sobre a mesa.

\begin{promptlabbox}{Escolher a ferramenta para a pergunta}
\textbf{Objetivo e contexto:} distinguir contagem, classificação e consulta. Considera três tarefas na biblioteca fictícia: contar as reservas de hoje, organizar mensagens por assunto e confirmar o estado de uma reserva.

\textbf{Instruções:} escreve a informação necessária e um modo de verificar o resultado de cada tarefa. Experimenta decidir se usarias uma regra, um modelo, uma consulta ou uma combinação.

\textbf{Resultado esperado e validação:} a contagem deve corresponder aos registos do dia. As categorias precisam de critérios e casos novos para as verificar. Para confirmar uma reserva, é preciso consultar o registo certo e interpretar o estado.

\textbf{Falhas possíveis:} contar mensagens em vez de reservas, presumir que toda a pergunta requer IA ou inventar um estado que o registo não mostra.
\end{promptlabbox}

A Ana ainda não sabe se pode sair de casa para levantar o livro. No próximo capítulo, vamos escrever uma resposta que lhe permita decidir o que fazer com a informação disponível.
